% Bagian: computability
% Bab: recursive-functions
% Subbagian: notations-pr-functions

\documentclass[../../../include/open-logic-section]{subfiles}

\begin{document}

\olfileid[id]{cmp}{rec}{not}
\olsection{Notasi Rekursi Primitif}

Salah satu keuntungan dari deskripsi induktif yang saksama tentang fungsi
rekursif primitif ialah kita dapat mendeskripsikannya secara sistematis.
Sebagai contoh, kita dapat memberikan suatu ``notasi'' kepada setiap fungsi
semacam itu sebagai berikut. Gunakan simbol $\Zero$, $\Succ$, dan
$\Proj{n}{i}$ untuk fungsi nol, penerus, dan proyeksi. Sekarang andaikan $h$
didefinisikan dengan komposisi dari fungsi berargumen $k$, yaitu~$f$, dan
fungsi-fungsi berargumen $n$, yaitu $g_0$, \dots,~$g_{k-1}$, serta kita telah
memberikan notasi $F$, $G_0$, \dots,~$G_{k-1}$ kepada fungsi-fungsi tersebut.
Dengan menggunakan simbol baru $\fn{Comp}_{k,n}$, kita dapat menyatakan fungsi
$h$ dengan $\fn{Comp}_{k,n}[F,G_0,\dots,G_{k-1}]$.

Untuk fungsi yang didefinisikan dengan rekursi primitif, kita dapat menggunakan
notasi yang serupa. Andaikan fungsi berargumen $(k+1)$, yaitu~$h$,
didefinisikan dengan rekursi primitif dari fungsi berargumen $k$, yaitu~$f$,
dan fungsi berargumen $(k+2)$, yaitu~$g$, sedangkan notasi yang diberikan
kepada $f$ dan~$g$ masing-masing adalah $F$ dan~$G$. Maka notasi yang diberikan
kepada~$h$ adalah $\fn{Rec}_k[F,G]$.

Ingat bahwa fungsi penjumlahan didefinisikan dengan rekursi primitif sebagai
\begin{align*}
  \Add(x_0, 0) & = \Proj{1}{0}(x_0) = x_0\\
  \Add(x_0, y+1) & = \Succ(\Proj{3}{2}(x_0, y, \Add(x_0, y))) = \Add(x_0, y) +1
\end{align*}
Di sini peran~$f$ dimainkan oleh $\Proj{1}{0}$, sedangkan peran~$g$ dimainkan
oleh $\Succ(\Proj{3}{2}(x_0, y, z))$. Fungsi terakhir diberi notasi
$\fn{Comp}_{1,3}[\Succ,\Proj{3}{2}]$ karena merupakan hasil pendefinisian
fungsi dengan komposisi dari fungsi berargumen~$1$, yaitu $\Succ$, dan fungsi
berargumen~$3$, yaitu $\Proj{3}{2}$. Dengan susunan ini, kita dapat menyatakan fungsi
penjumlahan dengan
\[
\fn{Rec}_1[\Proj{1}{0},\fn{Comp}_{1,3}[\Succ,\Proj{3}{2}]].
\]
Notasi-notasi ini kadang berguna, misalnya ketika mengenumerasi fungsi
rekursif primitif.

\begin{prob}
  Berikan notasi rekursif primitif lengkap untuk~$\Mult$.
\end{prob}

\end{document}
